\section{Piattaforma numerica}

\begin{frame}
\frametitle{Schema blocchi concettuale del procedimento}
\begin{center}
 \includegraphics[scale=0.5]{../../figure/flowchart/dataflow_platform}
 \end{center}

\end{frame}

%%%%%%%%%%%%%%%%%%%%%%%%%%%%%%%%%%%%%%%%%%%%%%%%%%%%%%%%%%%%%%%%%%%%%%%%%%%%%%%%%%%%%%%%%%%%%%%%%%%%%%%%
\begin{frame}
\frametitle{Modellazione dell'assembly di combustibile}
%%%%%%%%%%%%%%%%%%%%%%%%%%%%%%%%%%%%%%%%%%%%%%%%%%%%%%%%%%%%%%%%%%%%5
\begin{figure}[!htb]
\centering
\subfigure[Assembly di combustibile. \label{fassembly}]
{\includegraphics[width=0.48\textwidth]{../../figure/itex/fassembly}}
\centering
\subfigure[Elemento di combustibile. \label{unitcells}]
{\includegraphics[width=0.48\textwidth]{../../figure/itex/unitcells_wq}}
\caption{Modello di assembly di combustibile al fine di creare il database 
per il combustibile. }
\end{figure}
%%%%%%%%%%%%%%%%%%%%%%%%%%%%%%%%%%%%%%%%%%%%%%%%%%%%
\end{frame}

\begin{frame}
\frametitle{Modellazione dell'assembly di combustibile}
%%%%%%%%%%%%%%%%%%%%%%%%%%%%%%%%%%%%%%%%%%%%%%%%%%%%
\begin{table}[!htb]
\centering
 \begin{tabular}{lcc}
\toprule
Parametro & Unità & Valore \\
\midrule
Numero totale di esagoni & - & 169 \\
Numero di barre di combustibile & - & 126 \\
Numero di elementi di acciaio & - & 43 \\
Lato di un esagono & mm & 8.00208 \\
Apotema di un esagono & mm & 6.93 \\
Distanza tra barre di combustibile & mm & 13.86 \\
Raggio totale della barra di combustibile & mm & 5.25 \\
\bottomrule
\end{tabular}
\caption{Parametri di modellazione in riferimento all'assembly di combustibile 
in Figura \ref{fassembly}. \label{tabcrosection}}
\end{table}
%%%%%%%%%%%%%%%%%%%%%%%%%%%%%%%%%%%%%%%%%%%%%%%%%%%%%%%%%%

\end{frame}

\begin{frame}
\frametitle{Modellazione dell'assembly di combustibile}
%%%%%%%%%%%%%%%%%%%%%%%%%%%%%%%%%%%%%%%%%%%%%%%%%%%%%%5
\begin{table}[!htb]
\small
\centering
 \begin{tabular}{lcccccccc}
 \toprule
 Parametro & Unità & Valori &  &  &   & & \\
 \midrule
Temperatura del combustibile & [$^\circ$C] & 400.0 & 550.0 & 700.0 & 850.0 & 1000.0 
\\
\midrule
Temperatura del refrigerante & [$^\circ$C] & 380.0 & 440.0 & 600.0 & 750.0 &  \\
\midrule
Valori di bruciamento salvati & [Giorni] & 0.0 & 365.0 & 730.0 & 1095.0 & \\ 
\bottomrule
\end{tabular}
\caption{Dati di ingresso per la creazione del database combustibile.\label{daticompo}}
\end{table}
%%%%%%%%%%%%%%%%%%%%%%%%%%%%%%%%%%%%%%%%%%%%%%%%%%%%%
\end{frame}
%%%%%%%%%%%%%%%%%%%%%%%%%%%%%%%%%%%%%%%%%%%%%%%%%%%%%%%%%%%%%%%%%%%%%%%%%%%%%%%%%%%%%%%%%%%%%%%%%%%%%%%%
 \begin{frame}
 \frametitle{Modello implementato del reattore} 
 \begin{minipage}{0.48\linewidth}
 \begin{center}
 \includegraphics[scale=0.3]{../../figure/itex/hexagon}
 \end{center}
 \end{minipage}\hfill
 \begin{minipage}{0.48\linewidth}
 \begin{center}
 Sketch of the reactor core
 \begin{figure}
 \includegraphics[scale=0.3]{../../figure/itex/rectangle_wq}
 \end{figure}
 \end{center}
 \end{minipage}
   Reactor core fuel map 
 \end{frame}
 %%%%%%%%%%%%%%%%%%%%%%%%%%%%%%%%%%%%%%%%%%%%%%%%%%%%%%%%%%%%%%%%%%%%%%%%%%%%%%%%%%%%%%%%%%%%%%%%%%%%%%%%
 \begin{frame}
 %%%%%%%%%%%%%%%%%%%%%%%%%%%%%%%%
\begin{figure}[!htb]
\centering
\subfigure[Conducibilità termica mediata \label{tempeq}]
{\includegraphics[height=0.3\textheight]{./figure/itex/tempeq}}
\hspace{5mm}
\subfigure[Canale chiuso \label{canale}]
{\includegraphics[height=0.3\textheight]{./figure/itex/canale}}
\caption{Principali approssimazione considerate per la risoluzione termoidraulica di ciascun elemento del reattore.}
\end{figure}
%%%%%%%%%%%%%%%%%%%%%%%%%%%%%%%%
\begin{table}[!htb]
\centering
 \begin{tabular}{lccc}
\toprule
Parametro & Simbolo & Unità & Valore \\
\midrule
Densità & $\rho$ & kg/m$^3$ & 11340 \\
Conduttività termica & $\lambda$ & W/(m$\cdot$K) & 35 \\
Calore specifico & $c_p$ & J/(kg$\cdot$K) & 129 \\
Viscosità dinamica & $\mu$ & Pa$\cdot$s & 0.001844 \\
\bottomrule
\end{tabular}
\caption{Principali parametri termodinamici che caratterizzano il piombo 
liquido.}
\label{leadparam}
\end{table}
%%%%%%%%%%%%%%%%%%%%%%%%%%%%%%%%%%%%%%%%%%%%%%%%%%%%%%%
 
 
 
 
 
 \end{frame}
 
 
 
 
 \begin{frame}
 
 
 \end{frame}
